\documentclass[12pt,a4paper]{article}

% --- Encodage et langue ---
\usepackage[utf8]{inputenc}
\usepackage[T1]{fontenc}
\usepackage{lmodern}

% --- Maths ---
\usepackage{amsmath,amssymb,amsthm}

% --- Mise en page ---
\usepackage[top=2.5cm,bottom=2.5cm,left=2.8cm,right=2.8cm]{geometry}
\usepackage{parskip}

% --- Tableaux ---
\usepackage{booktabs}
\usepackage{array}

% --- Couleurs / liens ---
\usepackage{xcolor}
\usepackage{hyperref}
\hypersetup{colorlinks=true, linkcolor=blue!60!black, citecolor=blue!60!black}

% --- Environnements théorème ---
\newtheorem{theorem}{Théorème}[section]
\newtheorem{lemma}[theorem]{Lemme}
\theoremstyle{remark}
\newtheorem{remark}{Remarque}[section]

% --- Macros ---
\newcommand{\Sstar}{S^{*}}

% ======================================================
\begin{document}
% ======================================================

\begin{center}
  {\LARGE\bfseries Rapport de session}\\[6pt]
  {\large Extension multi-clause (Option 1) et limites du seuil $\Sstar$}\\[12pt]
  {\normalsize Validation empirique sur jeux de données synthétiques et réels}
\end{center}

\vspace{1em}
\hrule
\vspace{1.5em}

\tableofcontents

\newpage

% ======================================================
\section{Objectif de la session}
% ======================================================

Étendre la preuve de convergence existante (chaîne de Markov, 1 clause, cas
IDENTITY/NOT, \texttt{convergence\_sans\_bruit.tex}) à un réseau
\textbf{multi-clause}, selon la piste de partition fixe du domaine : la
partition est décidée avant l'apprentissage. L'espoir initial était que
chaque clause redevienne un sous-problème déjà couvert par la preuve
existante --- \textbf{mais attention, ça ne réduit pas l'entrée à 1 bit} :
une clause garde toutes ses $n$ features actives après masquage, seules les
features utilisées par le \emph{routage} sont gelées. Ce qui se réduit, dans
le meilleur cas, c'est le nombre de \textbf{littérales réellement
pertinentes} parmi ces features actives (le reste étant non pertinent, géré
par le même mécanisme d'annulation que dans le théorème IDENTITY/NOT). Et ce
n'est même pas garanti par construction : la section~\ref{sec:seuil-sstar}
montre des feuilles qui ont besoin de deux littérales pertinentes --- hors
du cadre du théorème actuel, qui ne couvre qu'une seule variable.

Fichiers créés pendant la session :
\begin{itemize}
  \item \texttt{tsetlin\_engine\_multiclause.h} --- moteur multi-clause
        (routeur, masquage, calibration, échantillonnage équilibré) ;
  \item \texttt{validation\_multiclause.cpp} --- XOR et DNF synthétiques ;
  \item \texttt{validation\_multiclause\_xor\_reel.cpp},
        \texttt{validation\_multiclause\_reels.cpp},
        \texttt{validation\_multiclause\_mnist.cpp} --- vrais jeux de
        données (\texttt{data/}).
\end{itemize}

Le moteur 1-clause d'origine, \texttt{tsetlin\_engine\_4case.h}, n'a
\textbf{jamais été modifié}.

% ======================================================
\section{Le routeur}
% ======================================================

Arbre de décision (CART simplifié) construit sur un batch $(x,y)$ séparé
(phase de découverte), puis \textbf{gelé} : pendant l'apprentissage il ne
prend que $x$ en entrée. Critère de split : impureté de Gini avec
\textbf{lookahead à 2 niveaux} (évalue le meilleur split atteignable une
fois la feature posée, pas seulement son gain immédiat) --- nécessaire pour
détecter les interactions de type XOR, où aucune feature seule n'apporte de
gain individuel.

$K$ est un budget fixé à l'avance, pas une conséquence découverte
automatiquement. Une tentative d'arrêter dynamiquement le découpage dès
qu'une feuille est représentable par une seule clause (test à la Find-S) a
été \textbf{abandonnée} : pour la classe des conjonctions, vérifier qu'une
solution exacte existe est aussi difficile que la trouver --- ce test
résolvait la formule cible en coulisses, et la clause n'avait alors plus
qu'à reproduire une réponse déjà connue. C'est une erreur méthodologique
réelle, détectée et corrigée en cours de session.

% ======================================================
\section{Résultats validés}
% ======================================================

Pour juger si une clause apporte une valeur réelle (et pas seulement le
découpage), chaque test compare le réseau complet à un \textbf{baseline
``routeur seul''} : le vote majoritaire de la feuille, sans aucun
apprentissage par clause.

\begin{center}
\begin{tabular}{@{}llccc@{}}
\toprule
Cible & $K$ & accuracy & baseline (routeur seul) & écart \\
\midrule
XOR synthétique & 2 & 100\% & 50.0\% & \textbf{+50.0\%} \\
XOR réel (40\% bruit train) & 2 & 100\% & 50.1\% & \textbf{+49.9\%} \\
DNF synthétique (insuffisant) & 2 & 75.0--82.6\% & 68.8\% & +6.2 à +13.8\% \\
DNF synthétique & 4 & 70.3--79.5\% & 73.0--75.0\% & $-4.7$ à $+4.5\%$ \\
AND2 réel (insuffisant) & 1 & 25.0\% & 75.0\% & $\mathbf{-50.0\%}$ \\
AND2 réel & 2 & 100\% & 75.0\% & \textbf{+25.0\%} \\
Maj3 réel (insuffisant) & 2 & 50.7\% & 75.5\% & $\mathbf{-24.8\%}$ \\
Maj3 réel & 5 & 100\% & 87.2\% & \textbf{+12.8\%} \\
MNIST (chiffre 0)\footnotemark & 2 & 93.9\% & 90.3\% & +3.6\% \\
MNIST (chiffre 0), $K$ croissant & 4--32 & 91.6--92.5\% & 90.3\% & $+1.3$ à $+2.2\%$ \\
\bottomrule
\end{tabular}
\end{center}
\footnotetext{Pour MNIST, le rappel sur la classe rare passe de $0\%$ ($K=2$
sans calibration) à $50.9\%$ ($K=2$, avec calibration $S/a$ +
échantillonnage équilibré) puis redescend à $14.1\%$ quand $K=32$ (dilution
du budget, cf. section~\ref{sec:seuil-sstar}). La colonne ``baseline'' pour
MNIST est le baseline trivial (toujours prédire la classe majoritaire), pas
le vote majoritaire par feuille --- non recalculé pour chaque $K$.}

XOR est la preuve la plus propre : dans chaque feuille, le vote majoritaire
est à $\sim$50\% (aucune majorité réelle), c'est la clause qui résout le bit
restant. Le routeur décide \emph{où regarder} ; la clause \emph{répond}.

Les lignes ``insuffisant'' (AND2 $K=1$, Maj3 $K=2$, DNF $K=2$) ne sont pas
des échecs séparés : elles partagent toutes la même cause, développée en
section~\ref{sec:seuil-sstar} --- au moins une feuille y apprend encore une
conjonction de 2 littéraux ou plus, ce que le seuil $\Sstar$ ne couvre pas à
40\% de bruit.

% ======================================================
\section{Bugs trouvés et corrigés}
% ======================================================

\begin{enumerate}
  \item \textbf{Masquage des features de routage.} Une feature utilisée
        pour atteindre une feuille devient constante dans cette feuille ; le
        mécanisme d'annulation $+S/-S$ qui exclut normalement les
        littérales non pertinentes ne s'applique plus. Sans correction, ces
        features dérivent vers l'inclusion à tort (cassait XOR). Corrigé en
        gelant ces automates à l'exclusion.
  \item \textbf{Feuilles ``constante à 0''.} Un automate figé exclu donne
        par défaut une sortie 1 (clause vide) ; une feuille qui doit
        toujours répondre 0 ne peut pas être représentée en excluant tout.
        Corrigé en forçant une littérale cohérente avec le chemin de
        routage (toujours violée) plutôt que de tout exclure.
  \item \textbf{Détection de pureté non tolérante au bruit.} Même défaut
        que Find-S (zéro exception tolérée), mais ici sur un simple
        comptage de fréquence --- donc corrigeable sans réintroduire la
        circularité. Remplacé par un seuil tolérant au bruit (\texttt{noiseTol}).
  \item \textbf{Routeur aveugle aux interactions (XOR réel).} Avec des
        features parasites, un critère à 1 niveau choisit une colonne de
        bruit au hasard plutôt que les vraies colonnes pertinentes (gain
        individuel nul pour XOR). Corrigé par le lookahead à 2 niveaux.
\end{enumerate}

% ======================================================
\section{Limites caractérisées (compromis réels, pas des bugs)}
% ======================================================

\subsection{Tension $S/a$ : vitesse de convergence vs. robustesse au bruit}

Sur une feuille dont la fonction restante est un OU (non représentable par
une seule clause), la clause converge vers le \textbf{ET le plus
restrictif} plutôt que vers le repli optimal (clause vide) --- pire que le
vote majoritaire. Mesuré directement sur le moteur 1-clause isolé : avec
$S \geq a$, convergence vers le pire cas (50\% sur une cible OU à 2
variables, contre 75\% pour la clause vide) ; avec $a > S$, convergence vers
le bon repli --- mais cela casse la convergence du cas bien représentable
(IDENTITY) à partir de $\sim$30--40\% de bruit.

\textbf{Correction générale apportée :} calibration de $(S,a)$ \textbf{par
feuille} (fonction \texttt{calibrateParams}), dérivée de la même analyse
$p^+/p^-$ que le théorème de convergence IDENTITY
(\texttt{convergence\_sans\_bruit.tex}). On mesure, pour chaque feuille, la
feature active la plus déséquilibrée et on calibre le rapport $S/a$ au
minimum nécessaire. Une bande de tolérance absorbe le bruit
d'échantillonnage normal pour ne pas sur-réagir.

\begin{remark}
\texttt{flip(p)} sature à $p \geq 1$ : il faut donc \textbf{abaisser $S$}
quand $a$ doit dominer, jamais pousser $a$ au-delà de 1 (sinon le rapport
réel retombe à 1:1, sans effet).
\end{remark}

\subsection{Déséquilibre de classe / features quasi-constantes (MNIST)}

Sur MNIST un-contre-tous (90\% non-cible / 10\% cible), le réseau ne prédit
jamais la classe rare (rappel $=0\%$). Cause : 353 des 784 pixels sont
quasi-toujours à 0 (fond d'image, indépendant du chiffre) --- même mécanisme
que ci-dessus mais dans les données brutes, pas dans le routage.

\textbf{Correction générale apportée :} un échantillonneur équilibré
(\texttt{BalancedSampler}, présente les exemples positifs/négatifs 50/50
pendant l'apprentissage, lu directement depuis les labels du jeu fourni,
sans ratio supposé à l'avance), combiné à la calibration ci-dessus.
Résultat : accuracy $90.3\% \to 93.9\%$, rappel $0\% \to 50.9\%$, avec les
mêmes paramètres de base partout ($S=1$, $a=0.3$, jamais réglés à la main
par jeu de données).

\subsection{Le sampling équilibré peut nuire (AND2, Maj3)}

Sur une cible AND2 ($x_1 \wedge x_2$, vraie 25\% du temps) avec 40\% de
bruit sur le label : dans le sous-ensemble étiqueté ``oui'', les faux
positifs sont \textbf{deux fois plus nombreux} que les vrais positifs,
simplement parce que les vrais négatifs sont trois fois plus fréquents au
départ. Forcer un tirage 50/50 amplifie l'exposition à ce bruit majoritaire
au lieu de le diluer. Résultat avec $K=1$ : accuracy $25\%$ (pire que le
hasard), clause vide.

C'est l'\textbf{inverse exact} du problème précédent : là, équilibrer
corrigeait un déséquilibre de \emph{features} ; ici, équilibrer un
déséquilibre de \emph{classe} amplifie le bruit sur la classe rare.
\textbf{Aucun correctif général trouvé cette session n'est gratuit.}

% ======================================================
\section{Découverte principale : le seuil $\Sstar$ ne couvre qu'une variable}
\label{sec:seuil-sstar}
% ======================================================

En cherchant la cause profonde de la section précédente, un test décisif a
été mené sur le moteur \textbf{1-clause d'origine, intact}, sans aucun ajout
de cette session :

\begin{center}
\begin{tabular}{@{}lc@{}}
\toprule
Bruit (cible $x_1 \wedge x_2$, 2 littéraux) & exactitude \\
\midrule
30\% & 100\% \\
35\% & 100\% \\
\textbf{40\%} & \textbf{0\% (accuracy 25\%, casse net)} \\
\bottomrule
\end{tabular}
\end{center}

Le moteur d'origine lui-même bascule entre 35\% et 40\% de bruit pour une
conjonction de \textbf{deux} variables --- alors que le théorème de
convergence IDENTITY ne couvre que les cibles à \textbf{une} variable. AND2
et Maj3 ont justement $\sim$40\% de bruit dans leurs fichiers réels : ils
tombent exactement dans la zone non couverte par la théorie actuelle.

Confirmation croisée : AND2 avec $K=2$ (qui découpe sur $x_2$, réduisant
chaque feuille à une seule variable pertinente) donne $100\%$ --- alors que
$K=1$ (qui force la clause à apprendre la conjonction des deux variables
directement) donne $25\%$. Le découpage n'aide pas par hasard : il
\textbf{réduit le problème au régime que la théorie couvre déjà}.

\begin{remark}
Ce n'est pas un bug --- c'est une lacune non caractérisée de la théorie
existante, qui explique d'un coup AND2, Maj3, et participe à la tension
$S/a$ de la section précédente.
\end{remark}

\subsection{Généralisation du seuil à $m$ littéraux}

Le théorème de convergence IDENTITY (avec bruit) ne couvre qu'une cible à
\textbf{une} variable. On généralise ici au cas d'une conjonction de $m$
littéraux positifs $x_1 \wedge \cdots \wedge x_m$, avec $x_i$ i.i.d.
Bernoulli($c{=}0.5$) et bruit de label symétrique $b$ (comme dans
\texttt{sampleLabeledBatch}). On reprend exactement le calcul $p^+/p^-$ du
théorème existant, appliqué à l'automate $v[1]$ (par symétrie, identique
pour $v[1..m]$).

Pour $x_1=0$, la conjonction est toujours fausse, donc $y_{\text{vrai}}=0$ ;
pour $x_1=1$, $y_{\text{vrai}}=1$ avec probabilité $\pi' = 2^{-(m-1)}$ (les
$m-1$ autres littéraux doivent aussi être à 1). En notant $a$ le paramètre
d'exclusion de l'engin (le $\alpha$ du théorème existant) :
\begin{align*}
p^+_{v[1]} &= P(x_1{=}0, y{=}0)\,S = \tfrac12(1-b)\,S, \\
p^-_{v[1]} &= P(x_1{=}0, y{=}1)\,a + P(x_1{=}1, y{=}0)\,S
            = \tfrac12 b\,a + \tfrac12\bigl[(1-\pi')(1-b)+\pi' b\bigr]S.
\end{align*}
La condition $\rho_{v[1]} = p^+_{v[1]}/p^-_{v[1]} > 1$ se simplifie (mêmes
étapes que pour le cas $m=1$) en :
\[
  S\,\pi'\,(1-2b) > b\,a
  \quad\Longleftrightarrow\quad
  \boxed{\;\Sstar(m,b) = \dfrac{a\,b\,2^{\,m-1}}{1-2b}\;}
\]
Le seuil critique de bruit $b^*(m)$ --- le point où $\Sstar(m,b^*)=1$, donc
où même $S=1$ (le maximum) ne suffit plus --- vaut :
\[
  b^*(m) = \dfrac{1}{2 + a\cdot 2^{\,m-1}}.
\]

\subsubsection*{Validation empirique (trois points indépendants, $a=0.3$)}

\begin{center}
\begin{tabular}{@{}cccl@{}}
\toprule
$m$ & $b^*$ prédit & Observé (moteur 1-clause intact) & Source \\
\midrule
1 ($a=1$) & 33.3\% & 33.3\% & déjà publié (Test A) \\
2 & 38.5\% & marche à 35\%, casse à 40\% & AND2 (cette session) \\
3 & \textbf{31.25\%} & marche à 30\%, transition à 31\%, cassé à 32\% & testé pour valider la formule \\
\bottomrule
\end{tabular}
\end{center}

Pour $m=3$, la prédiction (31.25\%) tombe exactement entre les deux points
de test les plus proches (30\% encore bon, 32\% déjà cassé), avec une
transition observée à 31\% (33\% d'exactitude, ni 100\% ni 0\%) --- la
signature attendue d'un seuil franchi de peu. La formule explique aussi
pourquoi $b^*$ \textbf{décroît} avec $m$ : chaque littéral supplémentaire
divise par deux la fréquence du signal positif, donc le bruit a
proportionnellement plus de poids face à un signal de plus en plus rare.

\begin{remark}
Cette formule donne, pour une cible et un niveau de bruit donnés, le $K$
minimal nécessaire \textbf{par calcul direct} (réduire chaque feuille jusqu'à
ce que son nombre de littéraux pertinents $m$ satisfasse $b < b^*(m)$) ---
au lieu du balayage empirique fait cette session pour AND2 et Maj3.
\end{remark}

\subsection*{Contre-exemple : MNIST ne suit pas ce principe}

Sur MNIST, augmenter $K$ ($2 \to 32$) \textbf{dégrade} systématiquement le
rappel ($50.9\% \to 14.1\%$), à l'opposé de AND2/Maj3. Hypothèse retenue
mais \textbf{non encore vérifiée} : dilution du budget (300 exemples de
découverte et $10^5$ pas d'apprentissage, répartis sur de plus en plus de
feuilles). Le principe ``découper jusqu'à une variable par feuille'' semble
spécifique aux cibles à structure logique propre (peu de littéraux
réellement responsables) ; MNIST n'a pas cette structure (motif diffus sur
toute l'image).

% ======================================================
\section{État actuel --- verdict}
% ======================================================

\begin{itemize}
  \item Le \textbf{mécanisme central} (partition figée + convergence par
        clause) est validé, y compris sur données réelles, avec une
        méthodologie (baseline routeur seul) qui prouve que les clauses
        font un travail réel.
  \item Le \textbf{système complet}, avec les correctifs généraux ajoutés
        (sampling équilibré, calibration $S/a$), n'est pas universellement
        robuste : chaque correctif aide certains jeux de données et nuit à
        d'autres. La calibration $S/a$ par feuille est la seule correction
        qui, jusqu'ici, n'a cassé aucun cas testé.
  \item La vraie limite identifiée n'est \textbf{pas} dans le code
        multi-clause, mais dans le seuil $\Sstar$ du moteur 1-clause
        d'origine, jamais caractérisé pour des cibles à plusieurs
        littéraux.
\end{itemize}

% ======================================================
\section{Recommandation pour la suite}
% ======================================================

Prioriser la généralisation théorique du seuil $\Sstar$ à une conjonction de
$m$ littéraux (étendre le théorème de convergence IDENTITY bruité), plutôt
que de continuer à corriger jeu de données par jeu de données. C'est la
pièce qui transforme les observations de cette session (AND2, Maj3, tension
$S/a$) en un résultat théorique unique et citable, et qui permettrait de
calculer directement, pour une cible donnée, le $K$ minimal nécessaire ---
au lieu de le découvrir par balayage empirique.

En secondaire : vérifier l'hypothèse de dilution de budget sur MNIST (batch
de découverte et budget d'entraînement proportionnels à $K$).

\end{document}
